\subsection{What the numbers do not show}\label{sec:threats}

\paragraph{Ground truth.} The keep-label bench takes its ground truth from the typeset source, not from a person: a word is a keep when the
mathematics of the \LaTeX{} file or the typewriter fonts of the PDF say so. That makes the labels cheap and repeatable, and it ties them to the
source. A word that the source does not mark but that a reader would still keep as written gets the label translate, and a rule that keeps it is
scored as a false positive. We have not measured how often this happens, so the precision in Table~\ref{tab:keep} is to be read with that in mind,
and the recall as the share of the labelled keeps that the rules find.

\paragraph{Controls and noise.} In the translation bench, 53 of the 312 probes are controls: lines that any engine should get right, with context or
without. A control shows whether a change moves something it should leave alone. In Table~\ref{tab:probes} the controls go from 35 to 37 correct
(66.0\% to 69.8\%) when the memory is added; the second run of the baseline sets the noise floor against which that difference is read. The sign test
counts only the probes on which the two runs disagree, 5 for A and 56 for M.

\paragraph{Cold start and warm cache.} The document reader is the one part whose cost depends on the state of the machine. On a cold start it compiles
its model, which can take tens of seconds; with a warm cache the same read is fast. That is the reason the lens does not wait for it
(Section~\ref{sec:overview}), and the reason the numbers of this section are warm numbers: the engine bench reports a warm median of
$0.15$\,s for a batch of three lines, and the committed results hold no tail latency and no cold-start figure. When successive answers are late, an
exponential back-off puts the service to rest for 1, 2, 4, up to 30\,s before the next request.

\paragraph{Baselines and bounds.} Every table compares a change with a baseline that we also run: the rules of one block before \code{CodeRegion} is
added, and condition A before the memory. Two bounds frame the numbers. The oracle that replays the references scores 312 of 312 under each of its
conditions, which is the upper bound of the probes; the small model is a floor, not a recommendation. We do not compare with another screen
translator: the benches measure what the screen adds to the translation problem, not how well any one engine translates.

\paragraph{Held-out sets.} Papers and cases each have a development set and a held-out set. The constants of the rules were fitted on the development
set, and the held-out set decides what we report; a gap between the two is the only measure of overfitting we can offer. On the cases the held-out set
scores higher than the development set (recall $1.00$ against $0.91$ once the region rule is added), the opposite of what overfitting would show. We read
that as a sign that the held-out cases are the easier ones, not as a result: 40 keep blocks are a small set.
